\documentclass[12pt]{article}

\usepackage[utf8]{inputenc}
\usepackage[T1]{fontenc}
\usepackage{amsmath,amssymb}
\usepackage{booktabs}
\usepackage{graphicx}
\usepackage[margin=2.5cm]{geometry}
\usepackage{xcolor}
\usepackage{hyperref}
\usepackage{enumitem}

\hypersetup{
  colorlinks=true,
  linkcolor=blue!60!black,
  citecolor=blue!60!black,
  urlcolor=blue!60!black,
  pdfauthor={David Alfyorov},
  pdftitle={Absence of Starobinsky inflation in the spectral action
    at conformal coupling}
}

\newcommand{\dd}{\mathrm{d}}
\newcommand{\erfi}{\operatorname{erfi}}
\newcommand{\order}{\mathcal{O}}

\title{Absence of Starobinsky inflation in the spectral action\\
  at conformal coupling}
\author{David Alfyorov\\[4pt]
\small\textit{davidich.alfyorov@gmail.com}}
\date{}

\begin{document}
\maketitle

\begin{abstract}
The spectral action $S = \mathrm{Tr}\,f(D^2/\Lambda^2)$ of the
standard Chamseddine--Connes spectral triple predicts conformal
coupling $\xi = 1/6$ for the Higgs field.  We show that this
prediction eliminates the $R^2$ term from the one-loop effective
action, removing the scalaron from the gravitational spectrum and
excluding Starobinsky inflation within the standard framework.
The $R^2$ coefficient $\alpha_R = 2(\xi - 1/6)^2$ vanishes
identically at $\xi = 1/6$, and $\xi = 1/6$ is an exact one-loop
fixed point of the renormalization group.  Eight mechanisms for
restoring the scalaron are examined; all fail within the standard
spectral triple.  The theory predicts exactly two gravitational
wave polarizations.  Detection of a scalar polarization would
falsify the standard spectral triple.
\end{abstract}

\noindent\textit{PACS:} 04.60.--m, 04.50.Kd, 98.80.Cq

\medskip

\noindent\textbf{1.}
The spectral action principle~\cite{CC:1996,CC:1997} derives the
bosonic action of gravity and the Standard Model from a single
operator trace $S = \mathrm{Tr}\,f(D^2/\Lambda^2)$, where $D$ is
the Dirac operator of an almost-commutative spectral triple
and~$\Lambda$ is the spectral cutoff.  The heat kernel expansion
of this trace yields the Einstein--Hilbert action, cosmological
constant, and curvature-squared corrections, with coefficients
determined by the Seeley--DeWitt coefficients
$a_{2k}$~\cite{Vassilevich:2003}.

The curvature-squared sector of the one-loop effective action
takes the form~\cite{BV:1987,BV:1990,CZ:2012}
\begin{equation}\label{eq:Gamma}
  \Gamma^{(1)} \supset \int \dd^4x\,\sqrt{g}\,
  \bigl[\alpha_C\, C_{\mu\nu\rho\sigma}C^{\mu\nu\rho\sigma}
  + \alpha_R(\xi)\, R^2\bigr]\,\ln\frac{\Lambda^2}{\mu^2}\,,
\end{equation}
where $\alpha_C = 13/120$ depends only on the Standard Model
particle content ($N_s = 4$, $N_D = 45/2$, $N_v = 12$) and is
independent of both~$\xi$ and the cutoff function~$f$, while the
$R^2$ coefficient is~\cite{Alfyorov:2026paper1}
\begin{equation}\label{eq:alphaR}
  \alpha_R(\xi) = 2\Bigl(\xi - \frac{1}{6}\Bigr)^{\!2}.
\end{equation}

\medskip

\noindent\textbf{2.}
In the Chamseddine--Connes spectral triple, the Higgs doublet
arises from the finite part of the Dirac operator, and the
curvature coupling $\xi R|H|^2$ is not a free parameter.  The
$a_4$ coefficient contains a common Yukawa-dependent factor
multiplying both the curvature coupling and the kinetic
term~$|\nabla H|^2$.  After canonical normalization, one
obtains~\cite{CCM:2006,CC:2010,CCS:2013,vS:2015,DLM:2014}
\begin{equation}\label{eq:xi}
  \xi = \frac{1}{6}\qquad\text{(conformal coupling)}.
\end{equation}
This result has been confirmed in five independent derivations
within the noncommutative geometry
framework~\cite{CCM:2006,CC:2010,CCS:2013,vS:2015,DLM:2014},
including Pati--Salam extensions, and is robust under
addition of BSM scalars from the finite Dirac
operator~\cite{BFS:2010,CCS:2013ps,vdD:2018}.

The one-loop beta function $\beta_\xi \propto (\xi - 1/6)$
vanishes at $\xi = 1/6$~\cite{ParkerToms:2009}, so conformal
coupling is an exact one-loop fixed point.  The full one-loop
RG equation for~$\xi$ in the SM with gravitational corrections
is~\cite{ParkerToms:2009,Buchbinder:1992}
\begin{equation}\label{eq:betaxi}
  \beta_\xi = \frac{1}{(4\pi)^2}\Bigl(\xi - \frac{1}{6}\Bigr)
  \bigl(6\lambda + 6y_t^2 - \tfrac{3}{2}g'^{\,2}
  - \tfrac{9}{2}g^2\bigr),
\end{equation}
where $\lambda$ is the Higgs self-coupling, $y_t$ the top
Yukawa coupling, and $g'$, $g$ the $\mathrm{U}(1)_Y$ and
$\mathrm{SU}(2)_L$ gauge couplings.  The prefactor
$(\xi - 1/6)$ ensures that $\xi = 1/6$ is a fixed point
for all values of the SM couplings.

\medskip

\noindent\textbf{3.}
The consequences are immediate.  At $\xi = 1/6$:
\begin{equation}\label{eq:alphaR-zero}
  \alpha_R\bigl(\xi = \tfrac{1}{6}\bigr) = 0\,.
\end{equation}
The $R^2$ term is absent from the one-loop effective action.
There is no scalaron.  Starobinsky inflation~\cite{Starobinsky:1980},
which requires a propagating scalar degree of freedom from the
$R^2$ term with mass $M_{\mathrm{inf}} \approx 1.28 \times
10^{-5}\,M_P$~\cite{Planck:2018}, is excluded in the standard
spectral action.

The scalar propagator denominator
\begin{equation}\label{eq:Pis}
  \Pi_s(z, \xi) = 1 + 6\Bigl(\xi - \frac{1}{6}\Bigr)^{\!2}\,z\,
  \hat{F}_2(z, \xi)
\end{equation}
reduces to $\Pi_s = 1$ identically at $\xi = 1/6$, for any
momentum~$z = k^2/\Lambda^2$ and any cutoff function~$f$.
The scalar gravitational mode does not propagate.  The
gravitational spectrum contains only the massless spin-2
graviton and the massive spin-2 fakeon with mass
$m_2 = \sqrt{z_0}\,\Lambda$~\cite{Alfyorov:2026paper1}.
The gravitational wave speed is $c_T = c$
exactly~\cite{Alfyorov:2026paper4}, compatible with the
GW170817 bound $|c_T - c|/c < 10^{-15}$~\cite{GW170817:2017}.

\medskip

\noindent\textbf{4.}
One may ask whether the scalaron can be restored within the
spectral action framework.  Table~\ref{tab:escape} examines
eight mechanisms; all fail.

\begin{table}[ht]
\centering
\caption{Mechanisms for restoring Starobinsky inflation
  in the spectral action, and the obstruction that closes each.}
\label{tab:escape}
\begin{tabular}{p{3.8cm}p{6.5cm}}
\toprule
Mechanism & Obstruction \\
\midrule
Sub-Planckian~$\Lambda$ &
  Conflicts with GUT interpretation of~$\Lambda$ \\
Large $\xi \sim 2 \times 10^4$ &
  Violates spectral-triple geometric boundary conditions \\
NCG $\sigma$-singlet~\cite{BFS:2010} &
  $\xi_\sigma = 1/6$ (conformal) by the same mechanism \\
Pati--Salam extension~\cite{CCS:2013ps} &
  All scalars acquire conformal coupling \\
Grand Symmetry~\cite{CCS:2013gs} &
  $\sigma$ generates Higgs mass, not scalaron \\
$N_s' \sim 4 \times 10^{10}$ BSM scalars &
  Requires BSM content incompatible with NCG \\
Modified cutoff~$f$ &
  $\alpha_R$ comes from $a_4$, independent of~$f$ \\
Two-loop corrections &
  No calculation exists; $\beta_\xi(\xi=1/6) = 0$ persists at one loop \\
\bottomrule
\end{tabular}
\end{table}

The only known paths are departures from the standard framework:
the zeta spectral action~\cite{KLV:2015} or the dilatonized
spectral action~\cite{CC:2006dilaton}.  These modify the
operator trace and cannot be obtained from the standard
$\mathrm{Tr}\,f(D^2/\Lambda^2)$.

\medskip

\noindent\textbf{5.}
The result provides a sharp falsification criterion.  Detection
of a scalar gravitational wave polarization (e.g., a
longitudinal breathing mode in the output of LIGO/Virgo/KAGRA
or LISA) would imply $\xi \neq 1/6$ and falsify the standard
Chamseddine--Connes spectral triple.  Current constraints on
scalar polarizations from binary pulsar timing~\cite{Stairs:2003}
and GW170814~\cite{GW170814:2017} are consistent with zero
scalar content.

At $\xi = 1/6$, the ratio $c_1/c_2 = -1/3$ in the
$\{R^2, R_{\mu\nu}^2\}$ Stelle basis is frozen under the
matter-only one-loop RG flow:
$\beta_{c_1}/\beta_{c_2} = -1/3$
identically~\cite{Alfyorov:2026paper1}.  This ratio
discriminates SCT from asymptotic safety ($c_1/c_2 \approx
-1.09$, truncation-dependent~\cite{BMS:2009}) and from the
perturbative asymptotic freedom branch
($c_1/c_2 \approx -0.326$~\cite{Shapiro:2004}).

\medskip

\noindent\textbf{6.}
To summarize: the standard spectral action predicts conformal
coupling $\xi = 1/6$, which is a one-loop RG fixed point.
At $\xi = 1/6$, the $R^2$ coefficient vanishes, the scalar
gravitational mode decouples, and Starobinsky inflation is
excluded.  Eight mechanisms for restoring the scalaron fail
within the standard framework.  The prediction of exactly two
gravitational wave polarizations is testable by current and
planned detectors.

\bigskip

I am grateful to Igor Shnyukov for collaboration and discussions.

\begin{thebibliography}{99}

\bibitem{CC:1996}
A.\,H.~Chamseddine, A.~Connes, Commun.\ Math.\ Phys.\
\textbf{186}, 731 (1997); hep-th/9606001.

\bibitem{CC:1997}
A.\,H.~Chamseddine, A.~Connes, J.\ Math.\ Phys.\
\textbf{47}, 063504 (2006); hep-th/0512169.

\bibitem{Vassilevich:2003}
D.\,V.~Vassilevich, Phys.\ Rept.\ \textbf{388}, 279 (2003);
hep-th/0306138.

\bibitem{BV:1987}
A.\,O.~Barvinsky, G.\,A.~Vilkovisky, Phys.\ Rept.\
\textbf{119}, 1 (1985).

\bibitem{BV:1990}
A.\,O.~Barvinsky, G.\,A.~Vilkovisky, Nucl.\ Phys.\ B
\textbf{333}, 471 (1990).

\bibitem{CZ:2012}
A.~Codello, O.~Zanusso, J.\ Math.\ Phys.\ \textbf{54},
013513 (2013); arXiv:1203.2034 [gr-qc].

\bibitem{Alfyorov:2026paper1}
D.~Alfyorov, ``Nonlocal one-loop form factors of the spectral
action with Standard Model content'' (2026);
doi:10.22541/au.177507193.33027854/v1.

\bibitem{CCM:2006}
A.\,H.~Chamseddine, A.~Connes, M.~Marcolli, Adv.\ Theor.\
Math.\ Phys.\ \textbf{11}, 991 (2007); hep-th/0610241.

\bibitem{CC:2010}
A.\,H.~Chamseddine, A.~Connes, J.\ Geom.\ Phys.\
\textbf{61}, 317 (2011); arXiv:0901.0577.

\bibitem{CCS:2013}
A.\,H.~Chamseddine, A.~Connes, W.\,D.~van Suijlekom,
J.\ Geom.\ Phys.\ \textbf{73}, 222 (2013);
arXiv:1304.7583.

\bibitem{vS:2015}
W.\,D.~van Suijlekom, \textit{Noncommutative Geometry and
Particle Physics} (Springer, 2015).

\bibitem{DLM:2014}
K.~van den Dungen, W.\,D.~van Suijlekom, Rev.\ Math.\ Phys.\
\textbf{24}, 1230004 (2012); arXiv:1204.0328.

\bibitem{BFS:2010}
T.\,D.~Brennan, S.\,E.~Flood, D.\,G.~Robbins, Eur.\ Phys.\
J.\ C \textbf{74}, 2862 (2014); arXiv:1401.1006.

\bibitem{CCS:2013ps}
A.\,H.~Chamseddine, A.~Connes, W.\,D.~van Suijlekom,
JHEP \textbf{11}, 132 (2013); arXiv:1304.8050.

\bibitem{vdD:2018}
K.~van den Dungen, J.\ Noncommut.\ Geom.\ \textbf{12}, 1155
(2018); arXiv:1711.07299.

\bibitem{CCS:2013gs}
A.\,H.~Chamseddine, A.~Connes, W.\,D.~van Suijlekom,
Fortschr.\ Phys.\ \textbf{62}, 797 (2014);
arXiv:1409.7574.

\bibitem{ParkerToms:2009}
L.~Parker, D.~Toms, \textit{Quantum Field Theory in Curved
Spacetime} (Cambridge, 2009).

\bibitem{Buchbinder:1992}
I.\,L.~Buchbinder, S.\,D.~Odintsov, I.\,L.~Shapiro,
\textit{Effective Action in Quantum Gravity} (IOP, 1992).

\bibitem{Starobinsky:1980}
A.\,A.~Starobinsky, Phys.\ Lett.\ B \textbf{91}, 99 (1980).

\bibitem{Planck:2018}
Y.~Akrami \textit{et al.} (Planck Collaboration), Astron.\
Astrophys.\ \textbf{641}, A10 (2020); arXiv:1807.06211.

\bibitem{Alfyorov:2026paper4}
D.~Alfyorov, ``Nonlinear field equations and FLRW cosmology
from the spectral action'' (2026); doi:10.5281/zenodo.19098027.

\bibitem{GW170817:2017}
B.\,P.~Abbott \textit{et al.} (LIGO/Virgo), Astrophys.\ J.\
Lett.\ \textbf{848}, L13 (2017); arXiv:1710.05834.

\bibitem{KLV:2015}
B.~Iochum, C.~Levy, D.~Vassilevich, Commun.\ Math.\ Phys.\
\textbf{340}, 1163 (2015); arXiv:1503.09054.

\bibitem{CC:2006dilaton}
A.\,H.~Chamseddine, A.~Connes, Phys.\ Rev.\ Lett.\
\textbf{99}, 071302 (2007); arXiv:0706.3690.

\bibitem{Stairs:2003}
I.\,H.~Stairs, Living Rev.\ Relativ.\ \textbf{6}, 5 (2003).

\bibitem{GW170814:2017}
B.\,P.~Abbott \textit{et al.} (LIGO/Virgo), Phys.\ Rev.\
Lett.\ \textbf{119}, 141101 (2017); arXiv:1709.09660.

\bibitem{BMS:2009}
D.~Benedetti, P.\,F.~Machado, F.~Saueressig, Nucl.\ Phys.\ B
\textbf{824}, 168 (2010); arXiv:0902.4630.

\bibitem{Shapiro:2004}
I.\,L.~Shapiro, Class.\ Quant.\ Grav.\ \textbf{25}, 103001
(2008); arXiv:0801.0216.

\bibitem{CPR:2008}
A.~Codello, R.~Percacci, C.~Rahmede, Ann.\ Phys.\
\textbf{324}, 414 (2009); arXiv:0805.2909.

\end{thebibliography}

\end{document}
